% ======================================================================
%  General Conclusion of The Actualization Theory
%  AT-MONO115 · Canonical Monograph (v2.0) · Zenodo-ready
%  Assembly only — no new physics, no new derivations.
% ======================================================================

\chapter*{General Conclusion}
\label{concl:conclusion}
\addcontentsline{toc}{chapter}{General Conclusion}

The Actualization Theory reconstructs physics from two primitives --- Difference and the
tensor reference $\eta$ --- through the canonical hierarchy

\[
\text{Difference} \;\longrightarrow\; \text{Actualization}
\;\longrightarrow\; \text{Inevitable Spectrum} \;\longrightarrow\; \text{Physics}.
\]

\section*{What has been established}

\begin{enumerate}
\item \textbf{The foundation.} Difference is the fundamental boundary: the irreducible
primitive whose count-conservation law is its definitional identity. Its two faces --- the
scalar $\rho$ and the tensor $\psi$ --- are the trace and traceless parts of the one
Difference. $\eta$ is the tensor reference against which the tensor face is read. The
minimal primitive set is $\{\text{Difference},\eta\}$.

\item \textbf{The process.} Actualization is the derived count-producing process of
Difference. Its convergence is the closure boundary $N=96$. The Closure Principle
characterizes the boundary; it does not derive the primitive.

\item \textbf{The spectrum.} The D96 spectrum is the inevitable output of the actualization
attractor: a unique, content-independent eigenspectrum carrying no choice and no primitive
status.

\item \textbf{The physics.} Quantum mechanics ($|\psi|^2=\rho$, measurement as
actualization), gravity and spacetime (emergent, native dynamics), the standard model
(derived reads of the D96 moments), and cosmology (a spectrum readout) are all derived from
the canonical foundation.

\textbf{Origin of Physical Units.} The D96 spectrum yields dimensionless ratios, moments,
and angles. Dimensional observables are obtained only after calibration to measured
anchors: $v$ for weak/Higgs masses, $m_e$ for quark masses, and the standard conversion
factors for SI gravity quantities. Those anchors fix the unit system; they are a boundary,
not a new derivation.

\item \textbf{The boundaries.} The limits of the framework are documented honestly:
Difference as the ontological boundary, the unresolved status of $\psi$, the numerical
boundary $\pi$, the Bekenstein $2\pi$ factor, and the hosted standard-model dynamics.
These are limits, not failures.

\item \textbf{The frontier.} There is no open derivation frontier. The frontier is the
experimental validation of the derived predictions: $S,T,U$, $a_e$, $0\nu\beta\beta$, and
the pre-registered predictions P1, P2, P3.
\end{enumerate}

\section*{What has not been claimed}

The monograph does not claim: a derivation of the hosted standard-model Lagrangian; a
derivation of the exact Bekenstein $1/4$ coefficient (it requires the imported $2\pi$
factor); a resolution of the ultimate ontological status of $\psi$; or any result beyond
the accepted derivations. The universality research program --- the operator basis across
organized domains, the lock law, the organization structure --- is treated as future work,
outside the core physics derivation.

\section*{The scientific status}

The Actualization Theory is falsifiable: it specifies explicit, pre-registered observations
whose absence or contradiction would force a revision. It claims exactly its derivations,
discloses its boundaries, and separates derivation from experimental validation. Its core
claim is that the physics of the standard model, gravity, and cosmology can be reconstructed
from Difference, Actualization, and the inevitable spectrum --- a claim whose decisive
confirmation now rests with the experiments that follow.

\vspace{1em}
\noindent\textbf{The Actualization Theory:}
\textit{a reconstruction of physics from Difference, Actualization and Spectrum ---
derived, bounded, and awaiting its experiments.}
